\documentclass[12pt]{article}
\usepackage[utf8]{inputenc}
\usepackage[T1]{fontenc}
\usepackage{hyperref}
\title{Kok Borok Language Resources for K-12 Curriculum Development}
\author{Compiled by Claude Sonnet 5 (claude-sonnet-5) for Andr\'as Kornai}
\date{}
\begin{document}
\maketitle

\section{Introduction}
Kok Borok (Kokborok; ISO 639-3 trp) is the main indigenous language of Tripura, India, and is also spoken in Bangladesh. It has been an official language of Tripura since 1979. The 2011 census counted about 880,000 speakers in Tripura. It is taught in schools.

\section{Primary Educational Resources}
\begin{itemize}
\item Tripura Government, Directorate of Kokborok \& Official Mother Languages. Per the portal, the Directorate prepares Kokborok textbooks for Classes I--XII. SCERT prints and distributes those for Classes I--VIII, and the Directorate itself does so for Classes IX--XII: \url{https://kokborokoml.tripura.gov.in/sites/default/files/2024-12/Kokborok.pdf}
\item CBSE Kokborok syllabi for Classes IX--X (secondary) and XI--XII (senior secondary): \url{https://cbseacademic.nic.in/web_material/CurriculumMain26/Sec/KOKBOROK_Sec_2025-26.pdf}, \url{https://cbseacademic.nic.in/web_material/CurriculumMain26/SrSec/KOKBOROK_SrSec_2025-26.pdf}. Earlier years are at \url{https://www.cbseacademic.nic.in/web_material/CurriculumMain24/Sec/KOKBOROK_Sec_2023-24.pdf} and \url{https://www.cbseacademic.nic.in/web_material/CurriculumMain25/SrSec/KOKBOROK_SrSec_2024-25.pdf}.
\end{itemize}
Individual textbook titles were not seen in the search results.

\section{Linguistic and Reference Materials}
\begin{itemize}
\item Tripura University Kokborok syllabus: \url{https://tripurauniv.ac.in/site/images/pdf/syllabus/Kokborok.pdf}
\item Listing of an old Kokborok grammar (bookseller page; details not verified): \url{https://exoticindia.com/book/details/kokborok-grammar-old-and-rare-book-nav904}
\end{itemize}
The search results mention a dialectal thesaurus published by the Directorate, but I did not see its URL or details.

\section{Religious and Community Resources}
None were found in these searches.

\section{Digital and Audio Resources}
None were found in these searches.

\section{Curriculum Development Context}
Kokborok has established state-level curriculum support in Tripura, with textbooks for Classes I--XII and CBSE syllabi for Classes IX--XII. Gaps are in detail: textbook titles, digital copies and audio, which were not verified here. Use of Kokborok as a medium of instruction (rather than as a subject) was not examined.
\end{document}